\section{差の確かさと大きさ：検定}

  \subsection{目で見た差とばらつき}

    ここまでは, 一つの集まりの平均について考えてきた.
    手元の標本から平均を計算し, その平均が取り直すたびにどれくらい動くかを標準誤差で見積もって, 母平均のいそうな範囲を信頼区間として捉えたのだった.

    しかし, 実際にデータを見るときには, 二つの集まりを比べたい場面が多い.
    新しい指導法を試したクラスと従来のクラス, 薬を飲んだグループと飲んでいないグループ, あるいは原作のある映画とない映画.
    素朴には, 二つの集まりそれぞれの平均を計算して, その差を見ればよさそうだ.
    差が大きければ「違いがある」と思えるし, 小さければ「たいした違いはない」と受け取れる.

    たとえば, ある学校で二つのクラスの英語テストの平均点を比べたとしよう.
    A組の平均は72点, B組の平均は68点だったとする.
    4点の差がある.
    「4点も違うのだから, 教え方に違いがあったのではないか」と感じるかもしれない.

    だが, この4点という差をそのまま信じていいだろうか.
    前章までに使ってきた標準誤差で確かめてみよう.
    各クラスの人数がそれぞれ15人で, 点数の標準偏差はどちらも12点だったとする.
    15人はそれぞれのクラスの全員だが, 知りたいのは教え方の違いだ.
    そこで, この15人を, 同じ教え方で学ぶかもしれない生徒たちの中からたまたま集まった顔ぶれと見なす.
    別の生徒たちが同じ教え方で学んでいたら平均点も違っていたはずで, 前章で見たように, 標本を取り直せば平均は動く.
    二つの組の点数を, 前章までと同じく互いに独立とみなすと, それぞれの組の平均の動きは無関係に生じる.
    引き算をしているからといって, ばらつきまで差し引かれるわけではない.
    一方の平均が偶然高く出て, もう一方の平均が偶然低く出れば, その差はむしろ広がってしまう.
    そのため, 差のばらつきは, それぞれの平均のばらつき（分散）の足し算になる.

    二つの集まりの標本サイズを $n_A, n_B$, 不偏分散を $u_A^2, u_B^2$ とすると, 差の標準誤差は次のように計算できる.
    \[
      \mathrm{SE}_{\text{差}} = \sqrt{\frac{u_A^2}{n_A} + \frac{u_B^2}{n_B}}
    \]
    先ほどのテストの数字を入れてみると,
    \[
      \mathrm{SE}_{\text{差}} = \sqrt{\frac{12^2}{15} + \frac{12^2}{15}} = \sqrt{9.6 + 9.6} \approx 4.4\;\text{点}
    \]
    となる.
    偶然で見込まれるばらつきは, およそ4.4点だ.
    二つの組の本当の平均がまったく同じだったとしても, たまたま集まった生徒の点数次第で, 平均の差は4.4点ほど動きうる.

    前章と同じように差の95\%信頼区間を作ってみると, $4 \pm 2 \times 4.4$, つまりおよそ $-4.8$ 点から $12.8$ 点になる.
    区間の中にゼロが含まれており, マイナスの値まで入っている.
    本当の平均では, むしろB組のほうが高い可能性すら残されているわけだ.
    平均点だけを見れば「4点も違う」と思えた差も, 一人ひとりの点数が散らばっている中では, 偶然のばらつきにすっぽり埋もれてしまう.

  \subsection{「差がない」を出発点にする}

    二つの集まりに差があるのかどうかを調べたいとき, 「差がある」ことを直接確かめようとすると足場が定まらない.
    2点の差も「差がある」だし, 10点の差も0.1点の差も「差がある」だ.
    差がある状態は無数にあって, どこを基準にして計算を始めればよいかが決まらないからだ.

    一方で, 「差がない（差がゼロ）」という状態はただ一つに決まる.
    そこで, まず「差がない世界」を仮に設定することから始める.
    二つの集まりの母平均はまったく等しく, 手元に見えている差はすべて偶然のばらつきにすぎない, と仮定してみるのだ.
    その仮定のもとで, 「今回手元に出たほどの差が偶然起こる確率はどれくらいか」を計算する.
    もしその確率がごく小さければ, 「差がないという仮定に無理があるのではないか」と考える.
    このように差があるかどうかを判断する手続きを\textbf{統計的検定}と呼ぶ.

    \begin{defi}{帰無仮説と対立仮説}{hypotheses}
      検定において基準として設定する「差がない」という仮説を\textbf{帰無仮説}と呼び, $H_0$ と表す.
      これに対して, 確かめたい主張である「差がある」という仮説を\textbf{対立仮説}と呼び, $H_1$ と表す.
    \end{defi}

    帰無仮説のもとで手元の差の珍しさを測るために, 差をそのばらつき（差の標準誤差）で割って標準化する.

    \begin{defi}{検定統計量}{test-stat}
      二つの標本平均の差を $\bar{x}_A - \bar{x}_B$, その差の標準誤差を $\mathrm{SE}_{\text{差}}$ とするとき,
      \[
        t = \frac{\bar{x}_A - \bar{x}_B}{\mathrm{SE}_{\text{差}}}
      \]
      で計算される値を\textbf{検定統計量}と呼ぶ.
    \end{defi}

    統計量とは, どこまでを同じとみなすかを決めて, 残る違いだけを拾う数であった.
    この検定統計量もまさにその考え方で作られている.
    テストの点数か視聴回数かという単位の違いや人数の違いがあっても同じ基準で比べられるように, 差を差の標準誤差で割って, 「差が偶然のばらつきの何倍にあたるか」という珍しさだけを拾い上げた数だ.

    英語テストの例で計算してみると, 差が4点で標準誤差が4.4点だから, 検定統計量は $4 / 4.4 \approx 0.9$ になる.
    差がない世界で標本を何度も取り直せば, 中心極限定理により, この値は0のまわりに釣鐘型に散らばる\footnote{標本サイズが小さいときは$t$分布と呼ばれる分布を使うが, そこそこ大きければ正規分布とほとんど変わらない.}.
    この分布の中で, 「手元の値と同じか, それ以上に極端な値が出る確率」を求める.
    これが\textbf{$p$値}だ.

    \begin{defi}{$p$値}{p-value}
      帰無仮説が正しいと仮定したもとで, 手元のデータから計算された検定統計量と同じか, それ以上に極端な値が生じる確率を\textbf{$p$値}と呼ぶ.
    \end{defi}

    英語テストの検定統計量 $0.9$ を例にすると, 差がない世界で $t$ がとる値の分布と $p$値の関係は\cref{fig:ch05-pvalue-tails}の左のようになる.
    $p$値は, 左右の両裾の面積を合わせたものだ.

    \begin{figure}[htbp]
      \centering
      % 第5章: 差がないと仮定して取り直したときの検定統計量の分布を二枚並べ, 両裾を塗る
% 左: 各組15人の t=0.9 (両裾およそ0.37), 右: 各組75人の t=2.0 (両裾およそ0.04)
% 曲線は標準正規密度 exp(-x^2/2)/sqrt(2 pi). 数値は本文どおり.
\begin{tikzpicture}
  \begin{axis}[
    width=0.44\linewidth, height=3.6cm, scale only axis,
    axis x line=bottom, axis y line=none,
    xmin=-4.2, xmax=4.2, ymin=0, ymax=0.47,
    xtick={-0.9,0,0.9}, xticklabels={$-0.9$,$0$,$0.9$},
    tick align=outside, x tick label style={font=\footnotesize},
    title={各組15人：$t=0.9$}, title style={font=\footnotesize, yshift=-2pt},
    clip=false,
  ]
    \addplot[draw=none, fill=black!45, domain=-4.2:-0.9, samples=80] {exp(-x^2/2)/sqrt(2*pi)} \closedcycle;
    \addplot[draw=none, fill=black!45, domain=0.9:4.2, samples=80] {exp(-x^2/2)/sqrt(2*pi)} \closedcycle;
    \addplot[black, very thick, domain=-4.2:4.2, samples=160] {exp(-x^2/2)/sqrt(2*pi)};
    \node[font=\footnotesize] at (axis cs:0,0.44) {$p$値 $\approx 0.37$};
  \end{axis}
\end{tikzpicture}
\hfill
\begin{tikzpicture}
  \begin{axis}[
    width=0.44\linewidth, height=3.6cm, scale only axis,
    axis x line=bottom, axis y line=none,
    xmin=-4.2, xmax=4.2, ymin=0, ymax=0.47,
    xtick={-2,0,2}, xticklabels={$-2.0$,$0$,$2.0$},
    tick align=outside, x tick label style={font=\footnotesize},
    title={各組75人：$t=2.0$}, title style={font=\footnotesize, yshift=-2pt},
    clip=false,
  ]
    \addplot[draw=none, fill=black!45, domain=-4.2:-2, samples=80] {exp(-x^2/2)/sqrt(2*pi)} \closedcycle;
    \addplot[draw=none, fill=black!45, domain=2:4.2, samples=80] {exp(-x^2/2)/sqrt(2*pi)} \closedcycle;
    \addplot[black, very thick, domain=-4.2:4.2, samples=160] {exp(-x^2/2)/sqrt(2*pi)};
    \node[font=\footnotesize] at (axis cs:0,0.44) {$p$値 $\approx 0.04$};
  \end{axis}
\end{tikzpicture}

      \caption{差がないと仮定して標本を取り直したときに, $t$統計量がとる値の分布. 左は英語テストの各組15人の場合, 右は各組75人の場合で, 灰色に塗った両裾の面積を合わせたものが$p$値である. 左は$t$が$-0.9$以下または$0.9$以上になる確率で, およそ$0.37$である. 右は$t$が$-2.0$以下または$2.0$以上になる確率で, およそ$0.04$である. 手元の差はどちらも4点だが, 人数が増えると$t$が大きくなり, 灰色の部分は小さくなる.}
      \label{fig:ch05-pvalue-tails}
    \end{figure}

    差がない世界での分布は0のまわりの釣鐘型だと分かっていて, 形がわかれば, どこにどれくらい集まるかまでわかるのだった.
    英語テストの検定統計量 $0.9$ についても, $-0.9$ 以下と $0.9$ 以上の両裾に入る割合がこれで決まる.
    合わせるとおよそ $0.37$ で, 表を引いて確かめることもできる.
    差がない世界であっても, 3回に1回くらいはこれほどの差が偶然起きてしまうということだ.
    まったく珍しいことではない.

    注意したいのは, $p$値は「帰無仮説が正しい確率」ではない, という点だ.
    あくまで「帰無仮説が正しいと仮定したもとで」計算された確率であって, 仮説そのものの正しさを表す数字ではない.
    この本のはじめに振ったサイコロを思い出してほしい.
    細工のないサイコロを100回振って1の目が30回以上出る確率を求め, 「千五百回に一回あるかないかの珍しいことが起きたか, そうでなければサイコロに細工がある」と述べた, その数が $p$値だったのだ.
    サイコロでは1の目が出すぎていないかだけを確かめたかったので多いほうの裾だけを数えたが, 多すぎても少なすぎても差があるかどうかを知りたいなら, 両裾を足す.
    その数は細工のないサイコロだと仮定して計算したもので, サイコロに細工がある確率ではない.

  \subsection{あらかじめ決めておく基準}

    $p$値がどこまで小さければ「珍しい」と判断してよいのだろうか.
    結果を見てから「$p = 0.07$ だが, まあ珍しいと言っていいだろう」とその都度基準を動かせるようにしておくと, 自分に都合のよい判断ばかりしてしまう.
    そこで, どこまで小さければ珍しいとみなすかの境界線を, データを見る前にあらかじめ決めておく.

    \begin{defi}{有意水準と検定の判定}{significance-level}
      帰無仮説を退ける基準として事前に定めた確率を\textbf{有意水準}と呼び, $\alpha$ と表す.
      $p$値が有意水準を下回ったとき, 結果は\textbf{有意}であるといい, 帰無仮説を\textbf{棄却}して対立仮説を採択する.
      下回らなかったときは, 帰無仮説を棄却しない.
    \end{defi}

    世の中で広く使われているのは $\alpha = 0.05$（5\%）という基準だ.
    これは前章の信頼水準95\%の裏返しにあたるもので, 特別な数理的必然があるわけではなく, 慣習として定着しているにすぎない.
    有意水準を5\%に決めておくということは, 「差が本当にないときにも, 同じ手順で検定すれば20回に1回ほどは『差がある』と判定してしまう」という誤りをあらかじめ承知の上で使う, ということだ.

    このように, 統計的検定には二つの誤りが必ずつきまとう.
    一つは, 本当は差がないのに「差がある」と判定してしまう誤りで, これを\textbf{第1種の過誤}と呼ぶ.
    この誤りを犯す確率は, 有意水準 $\alpha$ を超えないように制御されている.
    もう一つは, 本当は差があるのに見逃して「差があるとは言えない」としてしまう誤りで, これを\textbf{第2種の過誤}と呼ぶ.

    先ほどの各組15人の英語テストでは, $p$値がおよそ $0.37$ だったので, 有意水準5\%を大きく上回り, 帰無仮説を棄却できない.
    ただし, 差がゼロだと証明されたわけではない.
    裁判でいえば証拠不十分で無罪になったようなもので, 手元の15人のデータからは, 差があると言い切るだけの根拠が得られなかったにすぎない.
    もし本当に教え方の違いによる4点の差が存在していたとしても, たった15人のデータではばらつきに埋もれてしまい, 第2種の過誤として見逃してしまっている可能性がある.
    見逃しの誤りは, データを集める人数を増やすことで減らすことができる.

  \subsection{差の存在, 大きさ, 原因}

    では, 調べる人数を増やしてみたらどうなるだろうか.

    先ほどと同じ英語テストで, 今度は各組の人数を15人から75人に増やしたとしよう.
    手元の平均点と標準偏差は先ほどとまったく同じで, A組が72点, B組が68点, 標準偏差はどちらも12点だったとする.
    差の標準誤差を計算し直してみると,
    \[
      \mathrm{SE}_{\text{差}} = \sqrt{\frac{12^2}{75} + \frac{12^2}{75}} = \sqrt{1.92 + 1.92} \approx 1.96\;\text{点}
    \]
    となる.
    人数が増えたことで標準誤差は1.96点まで縮んだ.
    検定統計量は $4 / 1.96 \approx 2.0$ となり, $p$値はおよそ $0.04$ になる (\cref{fig:ch05-pvalue-tails}の右).

    $p$値が有意水準の0.05を下回った.
    差がない世界ではめったに起こらない珍しいことが手元のデータで起きたことになり, 帰無仮説は棄却される.
    「差がある」と言える根拠が立ったわけだ.

    ところが, ここで差の95\%信頼区間を作ってみると, 様子が変わる.
    $4 \pm 2 \times 1.96$ を計算すると, 差の区間はおよそ $0.1$ 点から $7.9$ 点になる (\cref{fig:ch05-ci-compare}).
    区間は確かにゼロを外れており, おおむね区間がゼロを含まないことと検定で有意になることは一致している.
    だが, ゼロでないところまでは絞り込めても, その差がたったの0.1点なのか, それとも8点近くもあるのかまでは, この区間からは絞り込めていない.

    \begin{figure}[htbp]
      \centering
      % 第5章 図5-1: 各組15人と各組75人のときの, 平均点の差の95%信頼区間
% 数値は本文どおり: 15人 4 ± 2×4.4 = -4.8〜12.8, 75人 4 ± 2×1.96 = 0.08〜7.92
\begin{tikzpicture}[x=0.42cm, y=1cm, font=\footnotesize]
  % 軸
  \draw[-stealth] (-7,0) -- (16,0) node[right] {点};
  \foreach \t in {-5,0,5,10,15} {
    \draw (\t,0) -- (\t,-0.08) node[below] {$\t$};
  }
  % 差がゼロの線
  \draw[densely dashed] (0,0) -- (0,2.3);
  % 各組15人
  \draw[line width=2.2pt, black!55] (-4.8,1.6) -- (12.8,1.6);
  \draw[thick] (-4.8,1.45) -- (-4.8,1.75) (12.8,1.45) -- (12.8,1.75);
  \fill (4,1.6) circle (2.4pt);
  \node[above] at (-4.8,1.75) {$-4.8$};
  \node[above] at (12.8,1.75) {$12.8$};
  \node[left, align=right] at (-7.2,1.6) {各組15人};
  % 各組75人
  \draw[line width=2.2pt, black!55] (0.08,0.8) -- (7.92,0.8);
  \draw[thick] (0.08,0.65) -- (0.08,0.95) (7.92,0.65) -- (7.92,0.95);
  \fill (4,0.8) circle (2.4pt);
  \node[above] at (0.9,0.95) {$0.1$};
  \node[above] at (7.92,0.95) {$7.9$};
  \node[left, align=right] at (-7.2,0.8) {各組75人};
\end{tikzpicture}

      \caption{英語テストのA組とB組の平均点の差について, 95\%信頼区間を各組15人のときと各組75人のときで並べた図. 横軸は差 (点). 黒い点は手元の差4点で, 縦の破線は差がゼロの位置を示す. 平均点 (A組72点, B組68点) と標準偏差 (どちらも12点) は, 二つの場合で同じである.}
      \label{fig:ch05-ci-compare}
    \end{figure}

    有意という判定は, 「差がある」と言える根拠であって, 「差が大きい」という根拠ではない.
    人数を増やせば標準誤差はいくらでも小さくなるため, 手元の差が0.01点のようなごく小さなものであっても, 人数を集めさえすれば $p$値は0.05を下回りうる.
    逆に人数が少なければ, それなりに大きな差があっても有意にならないことがある.
    $p$値の大小を決めているのは, 差の大きさそのものだけでなく, 人数とばらつきの組み合わせだ.

    さらに考えてみよう.
    仮に十分なデータを集めて, 確かに4点の差があると分かったとしよう.
    では, それは「教え方の違いが原因だ」と言ってよいのだろうか.
    そうとは言えない.
    A組にはもともと英語が得意な生徒が多く集まっていたのかもしれないし, 授業が行われた時間帯がA組のほうが集中しやすい午前中だったのかもしれない.

    差がある根拠はあるのか, その差はどれくらいの大きさなのか, それは何が原因なのか.
    この三つは別のことで, 必要な根拠もそれぞれ違う.
    じゃんけんの記録でいえば, 3ポイントの差が11{,}567回の上では偶然とは片づけにくいと言えたことは差の存在の話で, その差を当てにしてパーを出し続けても勝てる割合が2ポイントしか上がらないことは, 差を使って実際に得られる分の大きさの話だった.

  \subsection{映画の視聴回数に差はあるか}

    ここまで扱ってきた映画80本のデータに戻ってみよう.
    これから比べる二つの集まりは, ばらつきが同じかどうか分からないので, 英語テストのときと同じく, 同じだとは仮定せずに計算する\footnote{この前提で行う検定をWelchの$t$検定と呼ぶ.}.

    80本の映画を, 既存の漫画や小説などを映画化した「原作あり群」（37本）と, オリジナル脚本の「原作なし群」（43本）に分けてみる.
    それぞれの集まりの平均と標準偏差は次のとおりだった.
    \begin{itemize}
      \item 原作あり群（37本）: 平均 $274.1$ 万回, 標準偏差 $260.6$ 万回
      \item 原作なし群（43本）: 平均 $153.7$ 万回, 標準偏差 $209.7$ 万回
    \end{itemize}
    手元の平均の差は, $274.1 - 153.7 = 120.4$ 万回だ.
    原作のある作品のほうが, 平均して120万回ほど多く見られているように見える.

    この差について, 差の標準誤差を計算してみる.
    二つの集まりの視聴回数は, ここでも互いに独立とみなす.
    \[
      \mathrm{SE}_{\text{差}} = \sqrt{\frac{260.6^2}{37} + \frac{209.7^2}{43}} \approx 53.5\;\text{万回}
    \]
    差の標準誤差はおよそ53.5万回だ.
    ここから差の95\%信頼区間を求めると, $120.4 \pm 2 \times 53.5$ より, およそ $13$ 万回から $227$ 万回となる.
    区間はゼロを含んでいない.

    次に, 検定を行ってみよう.
    帰無仮説を「原作の有無によって平均視聴回数に差はない」とする.
    検定統計量は,
    \[
      t = \frac{120.4}{53.5} \approx 2.25
    \]
    となる.
    差がない世界で同じように標本を取り直したとき, この $2.25$ と同じかそれ以上に極端な値が出る確率（$p$値）を計算すると, およそ $0.028$ になる.
    差がないとして標本を取り直せば, これほどの差が出るのは100回に3回ほどしかない.
    有意水準の5\%を下回るため, 帰無仮説は棄却され, 原作の有無によって平均視聴回数に差がある根拠はある, と言える.

    ただ, 95\%信頼区間を見れば, その差は13万回程度かもしれないし, 227万回に達するかもしれない.
    13万回と227万回ではまるで別物だが, 手元のデータからはそこまでしか絞り込めていない.

    さらに, 原作があることが原因で視聴回数が伸びたのかどうかも, この計算からは言えない.
    原作のある作品は, もともと知名度が高かったり, 製作費が多く投じられていたり, 宣伝が大規模に行われていたりする傾向がある.
    二つの集まりの平均を比べただけでは, 原作の力によって視聴回数が増えたのか, それとも製作費や宣伝の規模が大きかったから増えたのかを取り出すことはできないのだ.
    原因を確かめるためには, 原作の有無だけが異なっていて, 製作費や宣伝規模といったほかの条件がすべて揃っているような二つの集まりを比べる必要がある.

    差がある根拠はある.
    だが, 大きさは13万回から227万回という広い範囲でしか言えず, 原作の有無が原因だとは言えない.

  \begin{mybox}{差の大きさを一つの数で言う}
    検定の結果は「差があると言えるかどうか」を答えてくれるが, 差の大きさそのものは教えてくれない.
    そこで, 差の大きさを標準偏差の何倍にあたるかで表す指標として, \textbf{Cohenの$d$}が使われることがある.
    \[
      d = \frac{\bar{x}_1 - \bar{x}_2}{u_d}
    \]
    ここで $u_d$ は, 二つの群のばらつきを代表させた標準偏差だ.

    二つの群でばらつきが違えば, どちらの標準偏差で割るかで $d$ の値が変わってしまう.
    そこで, 両群のばらつきを一つにまとめて分母にする.
    分散なら, 独立なばらつきどうしを足すだけで, 合わさったばらつきになる.
    この章で差の標準誤差を作ったときに使ったのがこの性質で, 標準偏差のままではこうはいかない.
    だからまとめるのは分散で行い, 最後に平方根で元の単位に戻す：
    \[
      u_d = \sqrt{\frac{u_1^2 + u_2^2}{2}}
    \]

    先ほどの映画の例で計算すると, $u_d \approx 236.5$ 万回となり, $d = 120.4 / 236.5 \approx 0.51$ となる.
    平均の差は, 各群の中でのばらつきのおよそ半分に相当する, ということだ.
    目安として $0.2$ なら小さい, $0.5$ なら中程度, $0.8$ なら大きいとされることが多いが, これも有意水準の0.05と同じで, 絶対の根拠があるわけではない.
  \end{mybox}

  \begin{mybox}{ここまででわかったこと}
    四段階の三つ目「ばらつきを見積もる」.

    \textbf{この章で言えるようになったこと}: 原作の有無による平均視聴回数の差は $120.4$ 万回で, $p$値は約 $0.028$ と, 差がある根拠は得られた.
    ただ, 差の95\%信頼区間は約 $13$ 万回から $227$ 万回で, 大きさはそこまでしか絞れていない.

    \textbf{まだ言えないこと}: 原作があるから多く見られたのか, 製作費や宣伝の違いによるものかは, この差だけでは言えない.
    原作以外の条件がそろった二つの集まりは, どうすれば手に入るのだろうか.

    \textbf{前の章の見方がどう変わったか}: 一つの平均に区間をつけたのと同じばらつきの計算で, 二つの平均の差も確かめられる.
    差がないとすれば今回の結果がどれほど珍しいかを調べても, 差の区間が狭くなるわけではなかった.
  \end{mybox}
